% =============================================================
% Capítulo 1. Introducción
% =============================================================

\chapter{Introducción}

Magic: The Gathering es un juego de cartas coleccionables con más de tres décadas de historia y una de las comunidades competitivas más extensas dentro de su categoría \cite{guinness}. Cada partida enfrenta a dos jugadores que emplean un mazo construido de antemano, y una parte determinante del resultado queda fijada antes del inicio de la partida, durante la fase de construcción. La selección de las sesenta cartas que componen un mazo constituye, por tanto, una de las decisiones estratégicas centrales del juego.

El formato Standard, sobre el que se centra este trabajo, restringe las cartas legales a las expansiones más recientes y renueva su contenido cada pocos meses \cite{wotc_standard}. Esta rotación mantiene activo el metajuego (el conjunto de mazos y estrategias predominantes en cada momento): los arquetipos dominantes varían con frecuencia y un mazo competitivo en un momento dado puede quedar obsoleto tras la siguiente actualización. La construcción de mazos para Standard se enfrenta, en consecuencia, a un objetivo cambiante y no estacionario.

Desde una perspectiva computacional, la construcción de mazos se formula como un problema de optimización combinatoria sujeto a restricciones duras y caracterizado por una función objetivo de evaluación costosa \cite{garciasanchez2016, bjorke2017}. El espacio de búsqueda comprende del orden de cuatro mil cartas legales \cite{scryfall}, de las que deben seleccionarse exactamente sesenta respetando un máximo de cuatro copias por carta y una base de tierras coherente con los colores del mazo. El número de combinaciones válidas alcanza una magnitud que descarta la enumeración exhaustiva como estrategia de resolución, lo que justifica el recurso a métodos aproximados. A esta dificultad se añade que la calidad de un mazo no puede deducirse directamente de su composición, ya que depende de su comportamiento en partidas reales frente a otros mazos, comportamiento que solo se conoce mediante el juego efectivo o su simulación.

La pregunta de investigación que articula este trabajo es si una metaheurística puede abordar este problema con una calidad comparable a la de un jugador experto, sin disponer del conocimiento de dominio que este acumula con la práctica. El problema reúne varias características que lo convierten en un banco de pruebas idóneo para la computación evolutiva, una familia de técnicas aplicada con éxito a numerosos problemas del ámbito de los juegos \cite{yannakakis2018, risi2017}. El espacio de soluciones es amplio y discreto. Las restricciones son rígidas y deben satisfacerse en todo momento, no únicamente en la solución final. Adicionalmente, la función objetivo es costosa y estocástica: la evaluación de un mazo exige simular partidas, cada una de las cuales requiere varios minutos de cómputo y arroja un resultado con una componente aleatoria que ningún cálculo cerrado captura.

Este último aspecto confiere al trabajo su principal interés metodológico. La mayoría de los enfoques previos para la generación automática de mazos recurren a funciones heurísticas que puntúan la composición de un mazo sin llegar a jugarlo, atendiendo a indicadores como la curva de maná, la proporción de tierras o la sinergia entre cartas \cite{garciasanchez2018, miernik2022, stiegler2016}. Tales medidas resultan económicas, pero constituyen aproximaciones indirectas del rendimiento real. En el presente trabajo la evaluación se realiza mediante la simulación de partidas completas con un motor que aplica el reglamento del juego, lo que plantea un reto concreto: cómo hacer evolucionar una población cuando la medida del fitness de cada individuo conlleva un coste de varios minutos y presenta varianza. La respuesta a este reto, que combina un sistema de emparejamiento eficiente con operadores que respetan la estructura del problema, constituye el núcleo técnico de la memoria.

Existe asimismo una motivación derivada del propio dominio. Un mazo no es una agregación arbitraria de elementos de alto valor individual, sino un sistema con un plan de juego coherente, en el que las cartas se seleccionan en conjuntos de cuatro copias y la base de tierras se dimensiona para sostener los colores que el mazo requiere. Un algoritmo que ignore esta estructura puede hallar combinaciones que ganan partidas en un entorno cerrado pero que ningún jugador reconocería como un mazo viable. Lograr que la metaheurística produzca mazos verosímiles, y no meramente eficaces frente a su propio entorno de evaluación, constituye uno de los hilos conductores del trabajo.

En consecuencia, este Trabajo de Fin de Grado propone un algoritmo genético específicamente adaptado a la construcción de mazos, cuyos operadores reproducen la lógica con la que un jugador humano compone un mazo de competición y cuya función de fitness se evalúa mediante la simulación real de partidas. La propuesta se valida de forma experimental sobre el formato Standard y se analiza tanto en su rendimiento como en su comportamiento emergente.

\section{Objetivos}
\label{sec:objetivos}

El objetivo general de este trabajo es desarrollar un algoritmo genético capaz de generar mazos competitivos y estructuralmente realistas para el formato Standard de Magic: The Gathering, evaluándolos mediante la simulación real de partidas.

Este objetivo general se concreta en los siguientes objetivos específicos:

\begin{enumerate}
    \item Modelar las restricciones estructurales y operativas del formato Standard y formalizar una representación del mazo adecuada para su tratamiento mediante operadores evolutivos.
    \item Diseñar operadores genéticos que incorporen la lógica de construcción de mazos de competición, en particular la norma de las cuatro copias, de modo que el espacio explorado se corresponda con el de las configuraciones que construye un jugador humano.
    \item Validar el rendimiento competitivo de las soluciones generadas, automatizando su evaluación en un entorno de simulación con un coste computacional acotado.
    \item Formular una función de fitness que integre el rendimiento real en partidas con una métrica de calidad estructural sensible al arquetipo de cada mazo.
    \item Analizar el comportamiento emergente del sistema, esto es, los arquetipos resultantes, la convergencia cromática y los límites del enfoque, documentando tanto sus aciertos como sus limitaciones.
\end{enumerate}

\section{Estructura de la memoria}

El resto del documento se organiza en seis capítulos, seguidos de la bibliografía y un anexo. El capítulo~\ref{cap:marco} presenta el marco teórico: las reglas de Magic relevantes para el problema, los fundamentos de los algoritmos genéticos, el simulador empleado como motor de evaluación y el estado del arte en construcción automática de mazos. El capítulo~\ref{cap:metodologia} expone la metodología de trabajo seguida durante el proyecto. El capítulo~\ref{sec:arquitectura} describe la propuesta, desde la representación del mazo hasta la función de fitness y el sistema de torneo. El capítulo~\ref{cap:implementacion} detalla la realización del sistema en código. El capítulo~\ref{cap:desarrollo} recoge la evaluación experimental: la evolución del proyecto desde el primer prototipo, los experimentos que validan cada decisión de diseño y el análisis de la ejecución definitiva. El capítulo~\ref{cap:conclusiones} recoge las conclusiones, el grado de cumplimiento de los objetivos y las líneas de trabajo futuro. Cierran la memoria la bibliografía y un anexo sobre la relación del trabajo con los Objetivos de Desarrollo Sostenible de la Agenda 2030.
